% Definir os termos necessários para compreender o documento.

%\begin{itemize}[leftmargin=1.2cm]
%    \item Requisitos funcionais serão identificados por \textbf{RF001, RF002, ...}
%    \item Requisitos não-funcionais serão identificados por \textbf{NF001, NF002, ...}
%    \item Prioridades adotadas: \textbf{Essencial}, \textbf{Importante} e \textbf{Desejável}
%\end{itemize}

Este documento adota convenções tipográficas e de identificação com o objetivo de padronizar a apresentação dos requisitos e facilitar sua leitura, rastreabilidade e manutenção. Os requisitos funcionais serão identificados pelo prefixo \textbf{RF}, seguido de um identificador numérico sequencial, no formato \textbf{RF001, RF002, RF003, ...}. De forma análoga, os requisitos não funcionais serão identificados pelo prefixo \textbf{RNF}, seguido de um identificador numérico sequencial, no formato \textbf{RNF001, RNF002, RNF003, ...}. Os identificadores são únicos e não devem ser reutilizados, mesmo em caso de remoção de requisitos.

As prioridades dos requisitos serão classificadas como \textbf{Essencial}, \textbf{Importante} e \textbf{Desejável}, sendo atribuídas individualmente a cada requisito, não havendo herança automática entre diferentes níveis de especificação. Além disso, termos relevantes, nomes de entidades e identificadores poderão ser destacados em \textbf{negrito} com o objetivo de facilitar a leitura e a compreensão do documento.